\newpage

% Configuracion de la seccion
\chapter*{Lista de Simbolos}
\addcontentsline{toc}{chapter}{Lista de Simbolos}
\markboth{\hfill LISTA DE SIMBOLOS}{\hfill LISTA DE SIMBOLOS}

% =======================================
% Lista de simbolos
% =======================================
\begin{Nomencl}[2em]

\item[$B_i$] Compensaci\'on instalada en le barra i.
\item[${b_{iI}}^m        $] M\'aximo valor de compensaci\'on inductiva admisible en la barra i.
\item[$CPtrue        $] Conjunto Pareto Óptimo.
\item[$Ptrue        $] Frente Pareto Óptimo.
\item[$E_k        $] Tensi\'on en el nodo k.
\item[$FP        $] Frente Pareto calculado.
\item[${f_1}^m        $] M\'axima potencial reactiva total nominal que se puede instalar en el EPS.
\item[$f_3        $] Desviaci\'on m\'axima de voltaje.
\item[$f:\mathbb{R}^n \rightarrow \mathbb{R}        $] Funci\'on objetivo.
\item[$G_{km}        $] Conductancia de l\'inea entre los nodos k y m.
\item[$k        $] Nodo de la red el\'ectrica.
\item[$\textbf{k}        $] N\'umero de funciones objetivos de un MOP.
\item[$l_i        $] N\'umero de l\'ineas asociadas a la barra i.
\item[$e(x)        $] Conjunto de restricciones de un MOP.
\item[$m        $] Nodo de una red el\'ectrica.
\item[$\textbf{m}        $] N\'umero de restricciones de un MOP.
\item[$n        $] N\'umero de barras del sistema el\'ectrico de potencia.
\item[$\textbf{n}        $] N\'umero de variables de decisi\'on.
\item[$O        $] Ómicron, par\'ametro utilizado en el Ómicron ACO.
\item[$P^*        $] La mejor de las soluciones guardadas por los algor\'itmos MOACO.
\item[$P_k        $] Potencia activa neta en la barra k.
\item[$Q_k        $] Potencia reactiva neta en la barra k.
\item[$S_k        $] Potencia aparente neta en la barra k.
\item[${V_i}^*        $] Voltaje deseado en la barra i.
\item[$\Delta V        $] Desviaci\'on de tensi\'on.
\item[$w_j        $] Peso de ponderaci\'on.
\item[$\textbf{X}        $] Espacio de decisi\'on.
\item[$\textbf{x} \in \mathbb{R}^n$] Variable de control o independiente.
\item[$Y_{kk}$] Admitancia propia del nudo $k$.
\item[$y_{km}$] Admitancia de l\'inea entre los nudos $k$ y $m$.
\item[$\textbf{y}$] Vector objetivo.
\item[$||.||$] Norma.
\item[$\triangleright$] Domina o es no comparable con.
\item[$\sim$] Son no comparables.
\item[$\beta$] Par\'ametro que determina la influencia de la visibilidad en un ACO.
\item[$\eta_1, \eta_2, \eta_3$] Funciones de visibilidad especialmente dise\~nadas para los
    algor\'itmos MOACOs.
\item[$O$] Ómicron.
\item[$\tau        $] Matriz de feromona utilizada por ACO.
\item[$\theta_k        $] Fase de la tensi\'on $E_k$.
\item[$x^*        $] Soluci\'on global \'optima.
\item[$u,v,x        $] Vectores de decisi\'on.
\item[$J_i        $] Conjunto de vecinos ubicados a un salto del nodo.
\item[$\Delta \tau        $] Cantidad de feromonas a depositar.
\item[$\eta        $] Visibilidad.
\item[$\alpha        $] Importancia relativa de las feromonas.
\item[$Z        $] Factor de normalizaci\'on.
\item[$f_s        $] Factor de suavizamiento.
\item[$\tau_{min}        $] Valor m\'inimo impuesto a la cantidad de feromonas.
\item[$q_o        $] Probabilidad de utilizar la regla pseudo aleatoria.
\item[$gb        $] Mejor soluci\'on global.
\item[$b        $] Cantidad de colonias.
\item[$l        $] N\'umero de hormigas.
\item[$P        $] Poblaci\'on de soluciones no nominadas.
\item[$\eta        $] Funci\'on de visibilidad global
\item[$\eta_1        $] Funci\'on de visibilidad que otorga mayor deseabilidad a las soluciones con
    poca compensaci\'on.
\item[$w_1        $] Valor de ponderaci\'on para $\eta_1$.
\item[$w_3        $] Valor de ponderaci\'on para $\eta_3$.
\item[$f_b        $] Factor de proporcionalidad para el c\'alculo de $w_2$.
\item[$M_2        $] M\'etrica de distribuci\'on del frente pareto.
\item[$M_4        $] M\'etrica de error al frente pareto \'optimo.
\item[$SP        $] M\'etrica de spacing.
\item[$HV        $] M\'etrica hipervolumen.
\item[$H_1        $] Hip\'otesis alternativa de una prueba de hip\'otesis.
\item[$\alpha-error        $] Probabilidad de cometer error de Tipo I.
\item[$\beta-error        $] Probabilidad de cometer error de Tipo II.
\item[$n_i        $] Tama\~no de la i-\'esima muestra.
\item[$\bar{x}        $] Media muestral para los test estad\'isticos.
\item[${x(i)}        $] i-\'esimo valor de una muestra.
\item[$\sigma        $] Desviaci\'on est\'andar o Desviaci\'on t\'ipica.
\item[$\sigma^2        $] Varianza de una poblaci\'on.
\item[$\mu        $] Valor que se desea alcanzar en el T-test.
\item[$t_c        $] Valor $t$ critico del T-test.
\item[$d        $] Numero de muestas del test de Levene.
\item[$W,p        $] Estad\'istico del test Shapiro-Wilk.
\item[$Y_{(n)}        $] Muestra ordenada para el test Shapiro-Wilk.

\end{Nomencl}
